\documentclass[10pt,a4paper]{amsart}
\usepackage[margin=19mm]{geometry}
\usepackage{amsmath,amssymb,mathtools}
\usepackage[scr=rsfs,cal=euler]{mathalfa}
\usepackage{xfrac}
\usepackage[unicode,hidelinks,bookmarks=false]{hyperref}
\newcommand{\ie}{i.e.\ }
\newcommand{\cf}{cf.\ }
\newcommand{\resp}{respectively\ }
\newcommand{\Subsection}{\subsection}
\providecommand{\nobreakdash}{\nobreak\hbox{-}}
\setlength{\emergencystretch}{2em}
\multlinegap0pt
\usepackage{bm}
\usepackage{bbm}
\usepackage{dsfont}
\usepackage{xspace}
\usepackage{cancel}
\usepackage{mathtools}
\usepackage{amsthm}
\makeatletter
\@ifundefined{thm}{\newtheorem{thm}{Theorem}}{}
\makeatother
\newcommand\an{\text{ann}\,}
\newcommand\res{\text{res}}
\newcommand\sub{\subseteq}
\newcommand\vv{\,\vert\,}
\title[Depth and Associated Primes in Group Cohomology]{Depth and Associated Primes in Group Cohomology}
\author{James A. Schafer}
\date{Source: arXiv:2509.06895v1 (CC BY 4.0)}
\begin{document}
\maketitle

\noindent\emph{Source and licence.} Depth and Associated Primes in Group Cohomology. Version arXiv:2509.06895v1 (CC BY 4.0), distributed under the Creative Commons Attribution 4.0 International licence, as recorded in the arXiv licence metadata for this exact version and shown on the abstract page https://arxiv.org/abs/2509.06895v1. Full rights notice: licenses/arxiv-2509.06895-CC-BY-4.0.txt.\par\smallskip

\noindent\textit{Editorial scope.} Counted unit: the Thm whose source label is \texttt{None} in the article (source0.tex lines 265--265, source label \texttt{None}), delivered together with the article's own complete original proof of it (source0.tex lines 266--285, one \texttt{proof} environment of 20 lines). No mathematics is written, repaired or re-derived: statement and proof are the article's own text line for line. Required context retained uncounted: none. Every definition, display and label of those lines is retained unchanged. Omitted from this package and not redistributed: 1--264, 286--318. Nothing inside a retained excerpt is omitted or altered: every retained body is delivered exactly as the article's lines are written, so no omission receipt of any kind accompanies this package. Citations: the retained excerpts carry no citation command at all, so no bibliography entry of the article is transcribed and no reference is redistributed. Reference adaptation: none was needed and none was made; every reference inside the delivered unit resolves inside it, so no unresolved reference or citation is produced. Printed slips kept literal: no printed word, symbol, hypothesis or number of the counted statement or of its proof was altered. Layout and class: the article's own class is \texttt{amsart} with packages including amsmath, amssymb, amscd, xy; the wrapper is the lane's pinned \texttt{amsart} editorial wrapper at 10pt on A4, which restates the theorem-like environments the retained text needs and adds the article's own macro definitions that the retained text needs (\texttt{\textbackslash an}, \texttt{\textbackslash res}, \texttt{\textbackslash sub}, \texttt{\textbackslash vv}), with extra wrapper packages (bm, bbm, dsfont, xspace, cancel, mathtools). The delivered numbering is the unit's own. Assets: no retained span contains a figure environment, a graphics-inclusion command, a PDF-page inclusion, a table or a TikZ picture; no image file is copied into this package, the package declares no asset dependency and no image file is redistributed with it. Licence: version arXiv:2509.06895v1 of this article is distributed under the Creative Commons Attribution 4.0 International licence, as recorded in the arXiv licence metadata for this exact version and shown on the abstract page https://arxiv.org/abs/2509.06895v1. A full rights notice is delivered at licenses/arxiv-2509.06895-CC-BY-4.0.txt. Characters: every retained excerpt is pure ASCII, so the delivered engine, XeLaTeX with its default Latin Modern text font, reports no missing character.
\begin{thm} $\frak p^{C_GE}_E=\text{rad}\,(\an(\res\,y))$. \end{thm}

\begin{proof} This will follow from a sequence of observations and propositions.  Let $P=\frak p^{C_GE}_E$.  Recall $J=\an\,\res\,y\sub P$ with $P\in \text{Ass}\,H^*(C_GE)$ and $\dim J=\dim P$.
\medskip
\par (1) If $(R,M)$ is a local or *-local ring with $\dim R=0$ then $M^s=(0)$ for some $s>0$.
\begin{proof} $R$ is Artinian so $M^s=M^{s+1}=M\, M^s$ for some $s>0$ and therefore by Nakayama's Lemma, $M^s=(0)$.\end{proof}

\par (2) $(R,M)$ local or *-local and $J\sub M$, homogeneous in the *-local case, then $\dim R/J=0$ if and only if $M^s\sub J$ for some $s>0$.
\begin{proof} $(\bar R,\bar M)=(R/J, M/J)$ is local or *-local.  If $\dim\bar R=0$ then $\bar M^s=(\bar 0)$, i.e. $M^s\sub J$.  If $M^s\sub J$ then $V(M)={M}\sub V(J)\sub V(M^s)=\cup_sV(M)={M}$.  Hence $V(J)={M}$. Here $V(I)$ denotes the homogenous spectrum of a homogeneous ideal. Therefore $\dim J=\dim R/J=0$. \end{proof}

Now let $R=H^*(C_GE)$ and let $R_{(P)}$ denote the homogeneous localization of $R$, i.e., only invert the homogeneous elements not in $P$.  Let $\pi:R\to R_{(P)}$ denote the natural map.

\par (3) If $I$ is a homogeneous annihilator ideal contained in $P\sub R$, then $I=\pi^{-1}I_{(P)}$.
\begin{proof} If $S$ denotes the set of homogeneous elements not in $P$, then  
$$ \pi^{-1}I_{(P)}=\{a\in R\vv \text{there exists } s\in S\text{ with }s\,a\in I \}.$$ If $I=\an\,U\sub P$ and $s$ is homogeneous not in $P$, $a\in R$ and $s\,a\in I$, then for all $u\in U$, $(s\,a)u=s(a\,u)=0$ and hence $a\,u=0$ since $s$ is a non-zro divisor. That is $a\in I$ and $I=\pi^{-1}I_{(P)}$. \end{proof}

\par (4) $J=\an\,\res\,y\sub P$ and $J$, $P$ are homogeneous of same dimension. Every homogeneous prime of $R_{(P)}$ is $Q_{(P)}$ for some unique homogeneous prime $Q\sub P$ of $R$.  Therefore if $J_{(P)}\sub Q_{(P)}\sub P_{(P)}$ which implies
$$ J=\pi^{-1}J_{(P)}\sub Q=\pi^{-1} Q_{(P)}\sub P=\pi^{-1} P_{(P)}.$$
Hence $Q=P$, $Q_{(P)}=P_{(P)}$ and $\dim J_{(P)}=0$.  Therefore there exists $s\geq 1$ with $(P_{(P)})^s\sub J_{(P)}$.
Hence $$ P^s\sub \pi^{-1}(P^s)_{(P)}= \pi^{-1}(P_{(P)})^s\sub \pi^{-1}J_{(P)}=J $$
and $P\sub\text{rad}\,J$. Since the converse is obvious $P=\text{rad}\,(\an\,\res\,y)$.
\end{proof}

\end{document}
